\documentclass{article}

\usepackage{preamble}

\title{ELEC2134 Lecture Notes}
\author{Oskar Meszaros (z5636164)}
\date{September 2026}

\begin{document}

\maketitle
\tableofcontents
\newpage

\section{Topic 1: Transform Methods}
\subsection{Waveforms and Periodic Functions}
\subsubsection{Waveform Fundamentals}
\begin{callout}{definition}[Definition]
\label{definition:t1-d1}
\begin{itemize}
    \item A waveform describes the shape of a signal as a function of time.
    \item A signal is a function that varies with time.
    \item A unit step signal stays at 0 until a certain time, then instantly becomes 1.
    \item A ramp signal increases linearly with time, starting from 0.
\end{itemize}
\end{callout}

The period ($T$) is the time taken for a waveform to complete one full cycle and start repeating
\[
    \text{frequency } f=\frac{1}{T}\text{Hertz (Hz)}.
\]
Hertz is cycles per second. The units of period is time.

\begin{callout}{example}[Example]
\label{example:t1-e1}
Some common frequencies.

\begin{table}[H]
    \centering
    \label{tab:t1-t1}
    \begin{tabular}{c|c}
        Frequency ($f$) & Period ($f=\frac{1}{T}$) \\
        \hline
        \hline
        1 kHz & $1\times 10^{-3}s=1ms$\\
        1 MHz & $1\times 10^{-6}s=1\mu s$ \\
        1 GHz & $1\times 10^{-9}s=1ns$ \\
        1 THz & $1\times 10^{-12}s=1ps$
    \end{tabular}
\end{table}
\end{callout}

The peak amplitude of a wave is the distance between the $x$ axis and the highest point of the wave. The peak-to-peak amplitude is two times this amplitude.

\begin{callout}{definition}[Periodic Waveforms (or Functions)]
\label{definition:t1-d2}

    A periodic waveform or function is any waveform or function that satisfies
    \[
        x(t)=x(t+T),
    \]
    where $t$ is the time variable, and $T$ is the period. In general,
    \[
        x(t)=x(t+nT),\forall n \in\mathbb{Z}.
    \]
\end{callout}

\subsubsection{Sinusoid waveforms}
A sinusoidal waveform is a type of period waveform, and can be described using the following equation
\[
    x(t)=A\sin(2\pi ft),
\]
where $x(t)$ is the instantaneous amplitude, $A$ is the peak amplitude, and $f$ is the frequency in Hz. It can also be described in terms of period with
\[
    x(t)=A\sin \left( \frac{2\pi }{T}t \right),
\]
where $T$ is the period. This can be simplified in terms of angular frequency $\omega =2\pi f=\frac{2\pi }{T}$ radians/sec as
\[
    x(t)=A\sin(\omega t).
\]

\subsubsection{Other Types of Waveforms}
\textbf{Square Waveform}

A square waveform is used mostly in digital electronic circuits for clock and timing control signals. It has very steep vertical transitions and flat top and bottoms.

It has pulse width $\tau =\frac{T}{2}$ and frequency $f=\frac{1}{T}$.

\noindent\textbf{Rectangular Waveform}

A rectangular waveform is similar to a square wave, however the duration of the high level (pulse width) is shorter than half the period. This is described by the duty cycle, which is the amount of time a signal is ON compared to the total time of one cycle.

\noindent\textbf{Triangular Waveform}

A triangular waveform is a non-sinusoidal waveform that linearly oscillates between a positive and negative peak. They are usually comprised of many sine or cosine waves with different frequencies and amplitudes.

\noindent\textbf{Sawtooth Waveform}

A sawtooth waveform is a periodic waveform whose shape resembles the teeth of a sawblade. It is similar to a triangular waveform, however the decay is more similar to a square wave's.

\subsubsection{Harmonic Components}
A harmonic is a multiple of the fundamental angular frequency of a waveform. The fundamental frequency is the frequency as defined above, giving the harmonics to be
\[
    \omega _0=2\pi f_0=\frac{2\pi }{T}\quad\Rightarrow\quad w_n=nw_0,f_n=nf_0,
\]
where $n$ is a positive integer.

The average value of a waveform is the mean level of the signal over one complete cycle. This is the overall DC component of the waveform. For a periodic waveform $x(t)$ with period $T$, this is
\[
    V_{avg}=\frac{1}{T}\int_0^Tx(t)dt.
\]

\subsection{Fourier Series}
\subsubsection{Frequency Domain Representation}
The frequency domain representation of a waveform shows how a time-domain (i.e. amplitude as a function of time) signal is distributed across different frequencies, effective decomposing a signal into it's component frequencies.

\begin{callout}{example}[Example]
\label{example:t1-e2}
What does the phase-domain of the sinusoid $v(t)$ below look like?
\[
    v(t)=\frac{4}{\pi }\sin(\omega_0t) + \frac{4}{3\pi }\sin(3\omega_0t)+\frac{4}{5\pi }\sin(5\omega_0t).
\]

\textbf{Answer:}
This would be a plot with three peaks; one at $\omega_0$, one at $3\omega _0$, and one at $5\omega_0$. The first one has a height of $4/\pi $, the second $43\pi $, and the last $4/5\pi $.
\end{callout}

\subsubsection{Fourier Series}
This process of breaking down a signal is called Fourier decomposition, resulting in a Fourier series. This is a sum of sinusoidal components with different frequencies and amplitudes.
The main application of a Fourier series is to determine the harmonic components that are present in a complex periodic waveform.

Any periodic waveform with a period of $T$ seconds can be represented by an infinite series of harmonically related sinusoids. The representation is known as the trigonometric Fourier series.
\[
    x(t)=A_0+\sum_{i=1}^{\infty} A_n\cos(n\omega _0t)+B_n\sin(n\omega _0t).
\]
The fundamental frequency of $x(t)$ is $\omega _0=2\pi f_0=\frac{2\pi }{T}$ radians/sec. The integer multiples of $\omega_0$ are known as harmonic frequencies of $x(t)$. The Fourier coefficients $A_0$, $A_n$ and $B_n$ are calculated directly from the signal $x(t)$ using the following equations
\begin{align*}
    A_0 &= \frac{1}{T}\int_{-T/2}^{T/2}x(t)dt \\
    \omega _0&=2\pi f_0=\frac{2\pi }{r}\text{ radians/sec} \\
    A_n&=\frac{2}{T}\int_{-T/2}^{T/2}x(t)\cos(n\omega _nt)dt\\
    B_n&=\frac{2}{T}\int_{-T/2}^{T/2}x(t)\sin(n\omega _nt)dt.
\end{align*}

\noindent\textbf{Fourier Series of an Even Periodic Function}

A function is defined as even if $x(t)=x(-t)$. For these functions/waveforms, the expressions for the Fourier coefficients reduce to
\[
    A_0=\frac{2}{T}\int_0^{T/2}x(t)dt,\quad A_n=\frac{4}{T}\int_0^{T/2}x(t)\cos(n\omega _0t)dt,\quad B_n=0.
\]
The reason that the $B$ coefficients disappear is because the periodic waveform is even, so there are no terms from the odd sine function.

\noindent \textbf{Fourier Series of an Odd Periodic Function}

A function is defined as odd if $x(t)=-x(t)$. For these function/waveforms, the expressions reduce to
\[
    A_0=A_n=0 \text{ for } n>0,\quad B_n=\frac{4}{T}\int_0^{T/2}x(t)\sin(n\omega _0t)dt
\]

\subsubsection{Reconstructing a Waveform}
To reconstruct (or synthesise) a waveform/function from a set of Fourier coefficients, simply plug them back into the above expression for Fourier series along with their index.

\subsection{Alternative Trigonometric Form}
In the frequency domain, the amplitude of each frequency component is determined by $A_n$ and $B_n$, with contributions from both the sine and cosine terms. We would now like to express the signal in an alternative form in which the amplitude of each frequency components is obtained using a single cosine function with a phase angle
\[
    x(t)=K_0=\sum_{n=1}^{\infty} K_n\cos(n\omega_0t+\phi _n).
\]
This is the amplitude-phase representation. We are able to find the phase angle $\phi _n$ with
\[
    \phi _n=\tan^{-1}\left( -\frac{B_n}{A_n} \right)=-\tan^{-1}\left( \frac{B_n}{A_n} \right).
\]
With this phase angle, we are also able to express the Fourier series using a complex number with $K_n\angle\phi _n=A_n-jB_n$.

\subsection{Complex Form of the Fourier Series}
The cosine-sine and amplitude-phase Fourier series from above are all trigonometric Fourier series. Another alternative Fourier series representations uses exponential functions instead of trigonometric functions. This form is therefore called the complex form of the (or exponential) Fourier series, and is given by
\[
    x(t)=\sum_{n=-\infty}^{\infty} C_ne^{jn\omega _0t},\quad C_n=\frac{1}{T}\int_0^Tx(t)e^{-jn\omega _0t}dt,
\]
where $\omega _0=\frac{2\pi }{T}$. The DC component is then given by $C_0$, and $C_n$ can be found with
\[
    C_n=\frac{1}{2}(A_n-jB_n),\quad |C_{-n}|=|C_n|=\frac{1}{2}\sqrt{A_n^2+A_n^2}.
\]
\begin{itemize}
    \item For even symmetry: $C_n=\frac{2}{T}\int_0^{T/2}x(t)\cos(n\omega _0t)dt$
    \item For odd symmetry: $C_n=-\frac{2j}{T}\int_0^{T/2}x(t)\sin(n\omega _0t)dt$
\end{itemize}

\subsection{Effect of Shifting}
Shifting a periodic function along either the vertical or the horizontal axis will generally affect its Fourier coefficients. These are also affected by the presence of certain types of functional symmetries in $f(t)$, particularly even-, odd-, and half-wave symmetry because some coefficient sets vanish altogether.

Shifting a periodic function $f(t)$ along either the vertical or the horizontal axis affects its amplitude-phase Fourier series. It is apparent that shifting $f(t)$ up or down the vertical axis changes only the average value $K_0$, leaving the amplitude $K_n$ and phase angle $\phi _n$ of each harmonics component unaffected.

\subsection{Half-Wave Symmetry}
An important type of symmetry is the half-wave symmetry. A periodic function is half-wave symmetric if
\[
    x(t)=-x \left( t-\frac{T}{2} \right).
\]
Shifting a half-wave symmetric wave or function by half a period then inverting it yields the original function. If a function has half-wave symmetric, then both $A_n$ and $B_n$ are 0 for even values of $n$.

\subsubsection{Odd Half-Wave Symmetry}
A periodic signal has odd half-wave symmetry is both of the following are true:
\[
    x(t)=-x \left( t-\frac{T}{2} \right),\quad \text{and}\quad x(t)=-x(-t).
\]

For a function with odd half-wave symmetry, the coefficients are given by
\begin{align*}
    A_n &= \begin{cases}
        \frac{4}{T} \int_0^{T/2}x(t)\cos(n\omega _0t)dt & \text{for n odd} \\
        0 & \text{for n even}
    \end{cases} \\
        B_n&=\begin{cases}
            \frac{4}{T}\int_0^{T/2}x(t)\sin(n\omega _0t)dt & \text{for n odd} \\
            0 & \text{for n even}
        \end{cases}
\end{align*}

For odd half-wave symmetry
\begin{itemize}
    \item All cosine terms are 0 $A_n=0$.
    \item All even harmonics are 0.
    \item Only odd sine harmonics appear in the Fourier series.
\end{itemize}

\subsubsection{Even Half-Wave Symmetry}
A periodic signal has even half-wave symmetry is both of the following are true:
\[
    x(t)=-x \left( t-\frac{T}{2} \right),\quad \text{and}\quad x(t)=x(-t).
\]

For a function with even half-wave symmetry, the coefficients are given by
\begin{align*}
    A_n &= \begin{cases}
        \frac{4}{T} \int_0^{T/2}x(t)\cos(n\omega _0t)dt & \text{for n odd} \\
        0 & \text{for n even}
    \end{cases} \\
        B_n&=\begin{cases}
            \frac{4}{T}\int_0^{T/2}x(t)\sin(n\omega _0t)dt & \text{for n odd} \\
            0 & \text{for n even}
        \end{cases}
\end{align*}

For even half-wave symmetry
\begin{itemize}
    \item All sine terms are 0 $B_n=0$.
    \item All sine harmonics are 0.
    \item Only odd cosine harmonics appear in the Fourier series.
\end{itemize}

\subsection{Examples}
\subsubsection{Waveforms}
\begin{callout}{example}[Example]
\label{example:eg1-1}
Sum these two sine waves
\begin{align*}
    f(t) &= 10\sin(200t) \\
    g(t) &= 8\sin(200t)
\end{align*}

\textbf{Answer:}
These sine waves have the same period $2\pi f=200$, $f=\frac{200}{2\pi }Hz=31.83Hz$.
\end{callout}

\newpage
\appendix
\section{Filters}
\subsection{Types of Filters}
The following are idealised specifications for the most common filter functions

\noindent\textbf{Ideal Low-Pass Filter}

Passes low frequencies, blocks high frequencies.
\[
    |H(j\omega )|=\begin{cases}
        1, & 0\leq\omega \leq\omega _c \\
        0, & \omega >\omega _c
    \end{cases}
\]
\indent Passband: $0\leq\omega <\omega _c$, stopband: $\omega _c<\omega \leq\infty$.

\noindent \textbf{Ideal High-Pass Filter}

Passes high frequencies, blocks low frequencies.
\[
    |H(j\omega )|=\begin{cases}
        0, & 0\leq\omega <\omega _c \\
        1, & \omega >\omega _c
    \end{cases}
\]
\indent Passband: $\omega _c<\omega \leq\infty$, stopband: $0\leq\omega <\omega _c$.

\noindent \textbf{Ideal Band-Pass Filter}

Passes only frequencies between $\omega _L$ and $\omega _H$.
\[
    |H(j\omega )|=\begin{cases}
        1, & \omega _L\leq\omega \leq\omega _H \\
        0, & \text{otherwise}
    \end{cases}
\]
\indent Passband: $\omega _L\leq\omega \leq\omega _H$, stopbands: $\omega <\omega _L$ and $\omega >\omega _H$.

\noindent \textbf{Ideal Band-Stop (Band-Reject) Filter}

Rejects frequencies between $\omega _L$ and $\omega _H$.
\[
    |H(j\omega )|=\begin{cases}
        0, & \omega _L\leq\omega _H \\
        1, & \text{otherwise}
    \end{cases}
\]
\indent Passbands: $\omega <\omega _L$ and $\omega >\omega _H$, stopband: $\omega _L<\omega <\omega _H$.

\end{document}
